\section{The exact two-level exponent ledger}
\label{sec:exponent-ledger}

The natural regime is the zero-loss corner of a more general
conditional ledger.  The purpose of this section is bookkeeping:
none of the displayed exponent profiles is proved for the physical
coefficient.

\begin{definition}[Two-sided exponent profile]
\label{def:exponent-profile}
Assume \(A_X\ne0\) for all sufficiently large \(X\).  Let
\[
 Q_X=X^{\beta+o(1)},\qquad
 J_X=X^{1-\beta+o(1)},\qquad
 N_{0,X}=J_XQ_X^2.
\]
A two-sided profile \((\mu,\lambda,\eta_{\canc})\), with
\(\mu,\lambda,\eta_{\canc}\ge0\), means
\begin{align}
 M_X&=X^{-\mu+o(1)}N_{0,X},
 \label{eq:mass-profile}\\
 D_X&=X^{-\lambda+o(1)}
       \frac{N_{0,X}^2}{Q_X},
 \label{eq:energy-profile}\\
 \canc_X&=X^{-\eta_{\canc}+o(1)}.
 \label{eq:cancellation-profile}
\end{align}
The exponent \(\eta_{\canc}\) measures relative cancellation against
the actual post-bin mass.  A finite profile of this form forces
\(Z_X\ne0\) eventually.  If \(Z_X=0\) on a subsequence, then
\(\Flat_X=0\) there and every flatness upper bound is automatic; no
finite \(\eta_{\canc}\) or \(\eta_Z\) is assigned on that zero-mode
subsequence.
\end{definition}

\begin{theorem}[Profile identities]
\label{thm:profile-identities}
Under \cref{def:exponent-profile},
\begin{align}
 \SeffX
 &=
 X^{\beta+\lambda-2\mu+o(1)},
 \label{eq:effective-profile}\\
 \Flat_X
 &=
 X^{\beta+\lambda-2\mu-2\eta_{\canc}+o(1)}.
 \label{eq:flatness-profile}
\end{align}
The TPC-32 bound
\[
 \Flat_X\le X^{\chi+o(1)}
\]
is certified by the profile exactly when
\begin{equation}
 \boxed{
 2\eta_{\canc}
 \ge
 \beta-\chi+\lambda-2\mu.}
\label{eq:relative-ledger}
\end{equation}
\end{theorem}

\begin{proof}
Divide the square of \eqref{eq:mass-profile} by
\eqref{eq:energy-profile}:
\[
 \SeffX
 =
 X^{-2\mu+\lambda+o(1)}Q_X,
\]
which is \eqref{eq:effective-profile}.  Multiply by
\(\canc_X^2\) and use the exact factorization to obtain
\eqref{eq:flatness-profile}.  Comparing its exponent with \(\chi\)
gives \eqref{eq:relative-ledger}.
\end{proof}

\begin{proposition}[Feasible mass--energy wedge]
\label{prop:feasible-wedge}
Any two-sided profile compatible with the verified TPC-32 envelopes
satisfies
\begin{equation}
 \boxed{\mu\le\lambda\le2\mu.}
\label{eq:feasible-wedge}
\end{equation}
Consequently,
\[
 \beta-\mu
 \le
 \beta+\lambda-2\mu
 \le
 \beta,
\]
subject also to the trivial nonnegative effective-support exponent.
\end{proposition}

\begin{proof}
The lower mass-to-energy inequality and \(S_X\ll Q_X\) give
\[
 D_X\gg X^{-2\mu+o(1)}
 \frac{N_{0,X}^2}{Q_X}.
\]
Comparison with \eqref{eq:energy-profile} yields
\(\lambda\le2\mu\).  Conversely,
\[
 M_X\ge\frac{D_X}{L_X}
 \gg X^{-\lambda+o(1)}N_{0,X}
\]
by \eqref{eq:atom-envelope}.  Comparison with
\eqref{eq:mass-profile} yields \(\mu\le\lambda\).
The effective-support bounds follow from
\eqref{eq:effective-profile}.
\end{proof}

\subsection{Reconciliation with the absolute zero-mode ledger}

Define the signed absolute zero-mode exponent \(\eta_Z\) by
\begin{equation}
 |Z_X|=X^{-\eta_Z+o(1)}Q_X^2.
\label{eq:absolute-Z-profile}
\end{equation}
It may be negative if the zero mode is larger than \(Q_X^2\).

\begin{proposition}[Relative-to-absolute exponent dictionary]
\label{prop:relative-absolute-dictionary}
Under the two-sided profile,
\begin{equation}
 \boxed{
 \eta_Z
 =
 \eta_{\canc}+\mu-(1-\beta).}
\label{eq:eta-dictionary}
\end{equation}
At the critical TPC-32 endpoint
\[
 \chi=\frac1{400},
 \qquad
 \beta=\frac{267}{400},
\]
the relative ledger \eqref{eq:relative-ledger} is exactly
\begin{equation}
 \boxed{2\eta_Z\ge\lambda.}
\label{eq:absolute-ledger}
\end{equation}
\end{proposition}

\begin{proof}
From \(|Z_X|=\canc_XM_X\),
\[
 |Z_X|
 =
 X^{-\eta_{\canc}-\mu+o(1)}
 J_XQ_X^2
 =
 X^{-\eta_{\canc}-\mu+1-\beta+o(1)}Q_X^2,
\]
which proves \eqref{eq:eta-dictionary}.  At the critical endpoint,
\[
 \beta-\chi
 =
 \frac{266}{400}
 =
 2(1-\beta).
\]
Substitute
\(\eta_{\canc}=\eta_Z-\mu+(1-\beta)\) into
\eqref{eq:relative-ledger}; all mass terms cancel, leaving
\(2\eta_Z\ge\lambda\).
\end{proof}

\begin{corollary}[Natural corner]
\label{cor:natural-corner}
At \(\mu=\lambda=0\), the critical relative requirement is
\[
 \eta_{\canc}\ge\frac{133}{400}=1-\beta,
\]
equivalently
\[
 \canc_X\le\frac{X^{o(1)}}{J_X}.
\]
\end{corollary}

\begin{proof}
Insert \(\mu=\lambda=0\) and
\(\beta-\chi=266/400\) into
\eqref{eq:relative-ledger}.
\end{proof}

\begin{remark}[No duplicated gain]
The mass exponent \(\mu\) appears once in the effective-support level
and once, with the opposite bookkeeping role, in the conversion from
relative to absolute cancellation.  Equation
\eqref{eq:absolute-ledger} shows that it cancels exactly.  Assigning
the same mass gain or loss again in the TPC-82 absolute ledger would
double count it \citep{WangTPC82}.
\end{remark}

\begin{definition}[Independent physical-loss ledger]
\label{def:physical-loss-ledger}
Let \(\Lambda_{\mathrm{phys}}\) be the sum of fixed-power losses
actually incurred by the full physical interface and downstream
comparison, excluding any loss already represented in
\(\lambda-2\eta_Z\).  The route continues only under the strict rule
\[
 \boxed{\Lambda_{\mathrm{phys}}<\frac1{400}.}
\]
Equality or excess is a stop condition.  This ledger is independent
of \eqref{eq:absolute-ledger}; the same loss is never charged in
both.
\end{definition}
